\subsection{Smart Contracts}

Moving past the basic building blocks of blockchain—like blocks, hashing, and consensus—we need to look at smart contracts. These programs take blockchain from a simple database and turn it into a programmable machine that can run actual code.\\

The idea of a "Smart Contract" goes back to Nick Szabo, who famously compared them to a vending machine\cite{szabo1997smartcontracts}. In the blockchain world, a smart contract is just code that runs directly on the network. Historically, once you deployed a smart contract, it was impossible to change because of how blockchain works. This caused massive problems during the famous 2016 DAO hack \cite{daoattack2017}. Because the code was locked forever, developers could not patch the security bug to stop the thief. This disaster actually split Ethereum into two separate chains: Ethereum Classic (ETC) and Ethereum (ETH). The ETC side stuck to the "Code is Law" rule, leaving the stolen history alone and staying with Proof-of-Work. Meanwhile, the main ETH chain chose to do a hard fork to force the stolen money back.\\

Today, contract code is still technically permanent once deployed on-chain. However, developers now bypass this issue by using a "Proxy Contract" setup. If you find a bug in your original contract (Contract A), you deploy a fixed version (Contract B) and tell the Proxy contract to route all user traffic to the new version instead. It is a bit like updating software without changing the user's link. Also, while ETC still uses mining, ETH eventually switched over to a Proof-of-Stake system in 2022. To prevent another nightmare like the DAO hack, deploying contracts today usually requires safety features like upgradeable proxies, emergency kill-switches, and professional code audits.\\

Vitalik Buterin described smart contracts as cryptographic "boxes" that hold value and only open up when specific rules are met\cite{buterin2014ethereum}. They automatically handle complex rules, update system data, and move digital assets around whenever predefined conditions are satisfied\cite{yaga2018blockchain}. For a school certificate platform, smart contracts act as the backend logic. They control exactly how a degree is generated, tie it to a student's public wallet address, handle global verification checks, or revoke a certificate if an administrative error happens. Most of the time, these contracts are written in a programming language called Solidity\cite{perez2020smartcontracts}.